%% ECON 282E -- Session 5, Deck A1: What Perturbation Computes
%% Block A opens the second road out of Session 3's "two roads".  S3 and S4 put
%% values on a grid and iterated; this deck asks what it would mean to
%% differentiate the equilibrium conditions instead, and what licenses it.
%% Source: Quant_Macro Lec.9 (Perturbation and Projection), frames 1-8.
%% That section follows Fernandez-Villaverde & Guerron-Quintana, Lectures on
%% Solution Methods for Economists V, closely; it is rewritten here and cited.
%% Numbers from tools/figures/s05_numbers.json, keys "three_methods",
%% "pert_first_order".

%% ECON 282E -- Foundations of Macroeconomics (UC Riverside, Fall 2026)
%% Shared Beamer preamble for every deck in slides/S01 ... slides/S10.
%%
%% Usage (a deck lives in slides/S0X/ and is compiled from that folder):
%%   %% ECON 282E -- Foundations of Macroeconomics (UC Riverside, Fall 2026)
%% Shared Beamer preamble for every deck in slides/S01 ... slides/S10.
%%
%% Usage (a deck lives in slides/S0X/ and is compiled from that folder):
%%   %% ECON 282E -- Foundations of Macroeconomics (UC Riverside, Fall 2026)
%% Shared Beamer preamble for every deck in slides/S01 ... slides/S10.
%%
%% Usage (a deck lives in slides/S0X/ and is compiled from that folder):
%%   \input{../preamble/course_preamble.tex}
%%   \decktitle{1}{A1}{Who I Am \& Why This Course}{September 24, 2026}
%%   \begin{document}
%%   \makedecktitle
%%   ...
%%
%% Adapted from Courses/AI-ECON-2026/source/shared_preamble.tex (Metropolis theme,
%% concept/caution/example/takeaway boxes, \sectiondivider). Changes for this course:
%%   * course title block (\decktitle / \makedecktitle) instead of \makebeamertitle
%%   * \zh{...} typesets Chinese (book titles) under pdflatex via CJKutf8
%%   * researchbox for the "Research in the wild" frames
%%   * section pages off; TikZ libraries and the course map loaded here
%% Build with pdflatex only (two passes); tools/build_slides.py does this for you.

\documentclass[10pt,english,aspectratio=169]{beamer}

\usepackage[T1]{fontenc}
\usepackage{textcomp}
\usepackage[utf8]{inputenc}
\usepackage{babel}
\usepackage{CJKutf8}

\usepackage{amsmath, amssymb, amsthm, graphicx}
\usepackage{booktabs, multirow, tabularx}
\usepackage{tikz}
\usetikzlibrary{positioning,arrows.meta,shapes,calc,fit,backgrounds}
\usepackage{pgfplots}
\pgfplotsset{compat=1.16}
\usepackage[most]{tcolorbox}
\tcbuselibrary{skins, breakable}
\usepackage{listings}
\usepackage{xcolor}

\lstset{
  basicstyle=\ttfamily\small,
  keywordstyle=\color{blue!80!black}\bfseries,
  commentstyle=\color{green!50!black}\itshape,
  stringstyle=\color{orange!80!black},
  backgroundcolor=\color{gray!8},
  frame=single,
  rulecolor=\color{gray!40},
  breaklines=true,
  showstringspaces=false,
  numbers=none,
}

\usetheme{metropolis}
\usefonttheme{professionalfonts}
\metroset{sectionpage=none}

\definecolor{LightGray}{RGB}{230,230,230}
\definecolor{RedTitle}{RGB}{0,0,200}
\definecolor{VeryLightGray}{RGB}{245,245,245}
\definecolor{DarkText}{RGB}{0,0,0}
\definecolor{UCRBlue}{RGB}{0,45,106}
\definecolor{UCRGold}{RGB}{241,171,0}

% Stage colours of the course arc (used by the course map and elsewhere)
\colorlet{StageTrunk}{blue!12}
\colorlet{StageBranch}{cyan!14}
\colorlet{StageMachine}{gray!18}
\colorlet{StageWall}{red!14}
\colorlet{StageTools}{orange!20}
\colorlet{StageData}{green!16}
\colorlet{StageAgents}{violet!14}

\hypersetup{bookmarks=false,breaklinks=true,pdfborder={0 0 0},colorlinks=true,
  linkcolor=DarkText,urlcolor=RedTitle}

\setbeamercolor{background canvas}{bg=white}
\setbeamercolor{title}{fg=RedTitle}
\setbeamercolor{frametitle}{bg=VeryLightGray,fg=DarkText}
\setbeamercolor{normal text}{fg=DarkText}
\setbeamercolor{alerted text}{fg=RedTitle}
\setbeamersize{text margin left=5mm, text margin right=5mm}
\addtobeamertemplate{frametitle}{}{\vspace{-0.5em}}

\setbeamertemplate{navigation symbols}{}
\setbeamertemplate{footline}{%
  \hspace{5mm}{\usebeamerfont{page number in head/foot}\color{black!45}%
  ECON 282E \,$\cdot$\, \insertshortdeck}\hfill%
  \usebeamercolor[fg]{page number in head/foot}%
  \usebeamerfont{page number in head/foot}%
  \insertframenumber\,/\,\inserttotalframenumber\kern1em\vskip2pt%
}

% ---------------------------------------------------------------------------
% Chinese text under pdflatex (book titles etc.)
\newcommand{\zh}[1]{\begin{CJK*}{UTF8}{gbsn}#1\end{CJK*}}

% ---------------------------------------------------------------------------
% Deck title block
%   \decktitle{session}{deck id}{title}{date}
\newcommand{\insertshortdeck}{}
\newcommand{\decksession}{}
\newcommand{\deckid}{}
\newcommand{\decktitletext}{}
\newcommand{\deckdate}{}
\newcommand{\decktitle}[4]{%
  \renewcommand{\decksession}{#1}%
  \renewcommand{\deckid}{#2}%
  \renewcommand{\decktitletext}{#3}%
  \renewcommand{\deckdate}{#4}%
  \renewcommand{\insertshortdeck}{Session #1 \,$\cdot$\, Deck #1.#2}%
  \title{#3}%
  \author{Zhigang Feng}%
  \date{#4}%
}
\newcommand{\makedecktitle}{%
  \begin{frame}[plain,noframenumbering]
    \begin{tikzpicture}[remember picture,overlay]
      \fill[UCRBlue] (current page.north west) rectangle ([yshift=-0.55cm]current page.north east);
      \fill[UCRGold] ([yshift=-0.55cm]current page.north west) rectangle ([yshift=-0.62cm]current page.north east);
      \node[anchor=west, text=white, font=\small\bfseries] at ([xshift=5mm,yshift=-0.275cm]current page.north west)
        {ECON 282E \,$\cdot$\, Foundations of Macroeconomics};
      \node[anchor=east, text=white, font=\small] at ([xshift=-5mm,yshift=-0.275cm]current page.north east)
        {UC Riverside \,$\cdot$\, Fall 2026};
    \end{tikzpicture}
    \vspace*{1.1cm}
    {\usebeamerfont{title}\color{black!55}\large Session \decksession\ \,$\cdot$\, Deck \decksession.\deckid\par}
    \vspace{0.25cm}
    {\usebeamerfont{title}\color{RedTitle}\LARGE\bfseries \decktitletext\par}
    \vspace{0.7cm}
    {\large Zhigang Feng\par}
    \vspace{0.15cm}
    {\color{black!65}\deckdate\par}
    \vfill
    {\footnotesize\itshape\color{black!55}Build it on paper $\;\to\;$ solve it on the machine $\;\to\;$ push it to the frontier with AI\par}
  \end{frame}}

% ---------------------------------------------------------------------------
% Section divider (plain frame, not counted)
\newcommand{\sectiondivider}[2]{%
\begin{frame}[plain,noframenumbering]
\begin{center}
\vfill
{\Large\textbf{#1}\par}
\vspace{0.35cm}
{\normalsize #2\par}
\vfill
\end{center}
\end{frame}
}

% ---------------------------------------------------------------------------
% Boxes
\newtcolorbox{conceptbox}[1][]{colback=blue!3, colframe=RedTitle!55, boxrule=0.35mm,
  arc=1mm, left=2mm, right=2mm, top=1mm, bottom=1mm, #1}
\newtcolorbox{cautionbox}[1][]{colback=red!3, colframe=red!50!black, boxrule=0.35mm,
  arc=1mm, left=2mm, right=2mm, top=1mm, bottom=1mm, #1}
\newtcolorbox{examplebox}[1][]{colback=green!3, colframe=green!45!black, boxrule=0.35mm,
  arc=1mm, left=2mm, right=2mm, top=1mm, bottom=1mm, #1}
\newtcolorbox{takeawaybox}[1][]{colback=VeryLightGray, colframe=RedTitle!55, boxrule=0.35mm,
  arc=1mm, left=2mm, right=2mm, top=1mm, bottom=1mm, #1}
\newtcolorbox{researchbox}[1][]{enhanced, colback=violet!4, colframe=violet!60!black,
  boxrule=0.35mm, arc=1mm, left=2mm, right=2mm, top=1mm, bottom=1mm,
  fonttitle=\bfseries\small, title={Research in the wild}, #1}

\newcommand{\compacttables}{\renewcommand{\arraystretch}{1.15}}
\newcommand{\src}[1]{{\tiny\color{black!55}#1}}

% ---------------------------------------------------------------------------
\input{../preamble/course_map.tex}

%%   \decktitle{1}{A1}{Who I Am \& Why This Course}{September 24, 2026}
%%   \begin{document}
%%   \makedecktitle
%%   ...
%%
%% Adapted from Courses/AI-ECON-2026/source/shared_preamble.tex (Metropolis theme,
%% concept/caution/example/takeaway boxes, \sectiondivider). Changes for this course:
%%   * course title block (\decktitle / \makedecktitle) instead of \makebeamertitle
%%   * \zh{...} typesets Chinese (book titles) under pdflatex via CJKutf8
%%   * researchbox for the "Research in the wild" frames
%%   * section pages off; TikZ libraries and the course map loaded here
%% Build with pdflatex only (two passes); tools/build_slides.py does this for you.

\documentclass[10pt,english,aspectratio=169]{beamer}

\usepackage[T1]{fontenc}
\usepackage{textcomp}
\usepackage[utf8]{inputenc}
\usepackage{babel}
\usepackage{CJKutf8}

\usepackage{amsmath, amssymb, amsthm, graphicx}
\usepackage{booktabs, multirow, tabularx}
\usepackage{tikz}
\usetikzlibrary{positioning,arrows.meta,shapes,calc,fit,backgrounds}
\usepackage{pgfplots}
\pgfplotsset{compat=1.16}
\usepackage[most]{tcolorbox}
\tcbuselibrary{skins, breakable}
\usepackage{listings}
\usepackage{xcolor}

\lstset{
  basicstyle=\ttfamily\small,
  keywordstyle=\color{blue!80!black}\bfseries,
  commentstyle=\color{green!50!black}\itshape,
  stringstyle=\color{orange!80!black},
  backgroundcolor=\color{gray!8},
  frame=single,
  rulecolor=\color{gray!40},
  breaklines=true,
  showstringspaces=false,
  numbers=none,
}

\usetheme{metropolis}
\usefonttheme{professionalfonts}
\metroset{sectionpage=none}

\definecolor{LightGray}{RGB}{230,230,230}
\definecolor{RedTitle}{RGB}{0,0,200}
\definecolor{VeryLightGray}{RGB}{245,245,245}
\definecolor{DarkText}{RGB}{0,0,0}
\definecolor{UCRBlue}{RGB}{0,45,106}
\definecolor{UCRGold}{RGB}{241,171,0}

% Stage colours of the course arc (used by the course map and elsewhere)
\colorlet{StageTrunk}{blue!12}
\colorlet{StageBranch}{cyan!14}
\colorlet{StageMachine}{gray!18}
\colorlet{StageWall}{red!14}
\colorlet{StageTools}{orange!20}
\colorlet{StageData}{green!16}
\colorlet{StageAgents}{violet!14}

\hypersetup{bookmarks=false,breaklinks=true,pdfborder={0 0 0},colorlinks=true,
  linkcolor=DarkText,urlcolor=RedTitle}

\setbeamercolor{background canvas}{bg=white}
\setbeamercolor{title}{fg=RedTitle}
\setbeamercolor{frametitle}{bg=VeryLightGray,fg=DarkText}
\setbeamercolor{normal text}{fg=DarkText}
\setbeamercolor{alerted text}{fg=RedTitle}
\setbeamersize{text margin left=5mm, text margin right=5mm}
\addtobeamertemplate{frametitle}{}{\vspace{-0.5em}}

\setbeamertemplate{navigation symbols}{}
\setbeamertemplate{footline}{%
  \hspace{5mm}{\usebeamerfont{page number in head/foot}\color{black!45}%
  ECON 282E \,$\cdot$\, \insertshortdeck}\hfill%
  \usebeamercolor[fg]{page number in head/foot}%
  \usebeamerfont{page number in head/foot}%
  \insertframenumber\,/\,\inserttotalframenumber\kern1em\vskip2pt%
}

% ---------------------------------------------------------------------------
% Chinese text under pdflatex (book titles etc.)
\newcommand{\zh}[1]{\begin{CJK*}{UTF8}{gbsn}#1\end{CJK*}}

% ---------------------------------------------------------------------------
% Deck title block
%   \decktitle{session}{deck id}{title}{date}
\newcommand{\insertshortdeck}{}
\newcommand{\decksession}{}
\newcommand{\deckid}{}
\newcommand{\decktitletext}{}
\newcommand{\deckdate}{}
\newcommand{\decktitle}[4]{%
  \renewcommand{\decksession}{#1}%
  \renewcommand{\deckid}{#2}%
  \renewcommand{\decktitletext}{#3}%
  \renewcommand{\deckdate}{#4}%
  \renewcommand{\insertshortdeck}{Session #1 \,$\cdot$\, Deck #1.#2}%
  \title{#3}%
  \author{Zhigang Feng}%
  \date{#4}%
}
\newcommand{\makedecktitle}{%
  \begin{frame}[plain,noframenumbering]
    \begin{tikzpicture}[remember picture,overlay]
      \fill[UCRBlue] (current page.north west) rectangle ([yshift=-0.55cm]current page.north east);
      \fill[UCRGold] ([yshift=-0.55cm]current page.north west) rectangle ([yshift=-0.62cm]current page.north east);
      \node[anchor=west, text=white, font=\small\bfseries] at ([xshift=5mm,yshift=-0.275cm]current page.north west)
        {ECON 282E \,$\cdot$\, Foundations of Macroeconomics};
      \node[anchor=east, text=white, font=\small] at ([xshift=-5mm,yshift=-0.275cm]current page.north east)
        {UC Riverside \,$\cdot$\, Fall 2026};
    \end{tikzpicture}
    \vspace*{1.1cm}
    {\usebeamerfont{title}\color{black!55}\large Session \decksession\ \,$\cdot$\, Deck \decksession.\deckid\par}
    \vspace{0.25cm}
    {\usebeamerfont{title}\color{RedTitle}\LARGE\bfseries \decktitletext\par}
    \vspace{0.7cm}
    {\large Zhigang Feng\par}
    \vspace{0.15cm}
    {\color{black!65}\deckdate\par}
    \vfill
    {\footnotesize\itshape\color{black!55}Build it on paper $\;\to\;$ solve it on the machine $\;\to\;$ push it to the frontier with AI\par}
  \end{frame}}

% ---------------------------------------------------------------------------
% Section divider (plain frame, not counted)
\newcommand{\sectiondivider}[2]{%
\begin{frame}[plain,noframenumbering]
\begin{center}
\vfill
{\Large\textbf{#1}\par}
\vspace{0.35cm}
{\normalsize #2\par}
\vfill
\end{center}
\end{frame}
}

% ---------------------------------------------------------------------------
% Boxes
\newtcolorbox{conceptbox}[1][]{colback=blue!3, colframe=RedTitle!55, boxrule=0.35mm,
  arc=1mm, left=2mm, right=2mm, top=1mm, bottom=1mm, #1}
\newtcolorbox{cautionbox}[1][]{colback=red!3, colframe=red!50!black, boxrule=0.35mm,
  arc=1mm, left=2mm, right=2mm, top=1mm, bottom=1mm, #1}
\newtcolorbox{examplebox}[1][]{colback=green!3, colframe=green!45!black, boxrule=0.35mm,
  arc=1mm, left=2mm, right=2mm, top=1mm, bottom=1mm, #1}
\newtcolorbox{takeawaybox}[1][]{colback=VeryLightGray, colframe=RedTitle!55, boxrule=0.35mm,
  arc=1mm, left=2mm, right=2mm, top=1mm, bottom=1mm, #1}
\newtcolorbox{researchbox}[1][]{enhanced, colback=violet!4, colframe=violet!60!black,
  boxrule=0.35mm, arc=1mm, left=2mm, right=2mm, top=1mm, bottom=1mm,
  fonttitle=\bfseries\small, title={Research in the wild}, #1}

\newcommand{\compacttables}{\renewcommand{\arraystretch}{1.15}}
\newcommand{\src}[1]{{\tiny\color{black!55}#1}}

% ---------------------------------------------------------------------------
%% Course map: the recurring "you are here" slide for ECON 282E.
%% Loaded by course_preamble.tex. Pattern follows AI-ECON-2026 lec01_map.tex, but the map
%% shows the whole ten-session course rather than one lecture.
%%
%%   \CourseMap{s}{label}   s = 1..10 highlights session s and prints "you are here: label"
%%                          (e.g. \CourseMap{1}{1.A2}); earlier sessions get a check mark.
%%   \CourseMap{0}{}        full overview, nothing highlighted.
%%
%% Note: TikZ option lists are comma-split before TeX conditionals run, so every \ifnum
%% below sits outside [...] option lists.

\newcommand{\CourseMap}[2]{%
\begin{center}
\begin{tikzpicture}[
    sess/.style={rectangle, rounded corners=2pt, draw=black!45, minimum width=1.34cm,
        minimum height=1.12cm, text width=1.26cm, align=center, inner sep=1pt},
    stage/.style={font=\tiny\bfseries, text=black!70},
]
  \foreach \i/\d/\t/\c in {%
      1/Sep 24/Intro \\ Growth/StageTrunk,
      2/Oct 1/HA $\cdot$ OLG \\ Policy/StageBranch,
      3/Oct 8/Num.\ DP \\ Python/StageMachine,
      4/Oct 15/Numerics \\ I/StageMachine,
      5/Oct 22/Numerics \\ II/StageMachine,
      6/Oct 29/The wall \\ \textbf{Midterm}/StageWall,
      7/Nov 5/ML \\ Deep nets/StageTools,
      8/Nov 12/RL \\ HA $+$ DL/StageTools,
      9/Nov 19/Text \\ LLMs/StageData,
      10/Dec 3/Agents \\ \textbf{Projects}/StageAgents} {
    \node[sess, fill=\c] (n\i) at ({1.46*(\i-1)}, 0)
      {{\scriptsize\bfseries S\i}\\[-1pt]{\tiny\color{black!60}\d}\\[1pt]{\tiny \t}};
    \ifnum\i<#1
      \node[font=\scriptsize, text=green!45!black] at ([xshift=-2.5mm,yshift=-2.2mm]n\i.north east) {$\checkmark$};
    \fi
    \ifnum\i=#1
      \draw[RedTitle, very thick, rounded corners=2pt]
        ([xshift=-0.6mm,yshift=-0.6mm]n\i.south west) rectangle ([xshift=0.6mm,yshift=0.6mm]n\i.north east);
      % keep the label inside the picture at both ends of the row
      \ifnum\i<3
        \node[font=\scriptsize\bfseries, text=RedTitle, anchor=south west] at ([yshift=1.2mm]n\i.north west)
          {$\blacktriangledown$\,you are here: #2};
      \else\ifnum\i>8
        \node[font=\scriptsize\bfseries, text=RedTitle, anchor=south east] at ([yshift=1.2mm]n\i.north east)
          {you are here: #2\,$\blacktriangledown$};
      \else
        \node[font=\scriptsize\bfseries, text=RedTitle, anchor=south] at ([yshift=1.2mm]n\i.north)
          {you are here: #2\,$\blacktriangledown$};
      \fi\fi
    \fi
  }
  % stage band under the sessions
  \foreach \a/\b/\lab/\c in {1/1/trunk/StageTrunk, 2/2/branches/StageBranch,
      3/5/the machine $\cdot$ classical methods/StageMachine, 6/6/the wall/StageWall,
      7/8/new tools/StageTools, 9/9/new data/StageData, 10/10/colleagues/StageAgents} {
    \fill[\c] ([yshift=-1.5mm]n\a.south west) rectangle ([yshift=-4.6mm]n\b.south east);
    \path ([yshift=-1.5mm]n\a.south west) -- ([yshift=-4.6mm]n\b.south east)
      node[midway, stage] {\lab};
  }
  \node[font=\footnotesize\itshape, text=black!70] at ([yshift=-9.5mm]$(n1.south)!0.5!(n10.south)$)
    {Build it on paper $\;\to\;$ solve it on the machine $\;\to\;$ push it to the frontier with AI};
\end{tikzpicture}
\end{center}
}


%%   \decktitle{1}{A1}{Who I Am \& Why This Course}{September 24, 2026}
%%   \begin{document}
%%   \makedecktitle
%%   ...
%%
%% Adapted from Courses/AI-ECON-2026/source/shared_preamble.tex (Metropolis theme,
%% concept/caution/example/takeaway boxes, \sectiondivider). Changes for this course:
%%   * course title block (\decktitle / \makedecktitle) instead of \makebeamertitle
%%   * \zh{...} typesets Chinese (book titles) under pdflatex via CJKutf8
%%   * researchbox for the "Research in the wild" frames
%%   * section pages off; TikZ libraries and the course map loaded here
%% Build with pdflatex only (two passes); tools/build_slides.py does this for you.

\documentclass[10pt,english,aspectratio=169]{beamer}

\usepackage[T1]{fontenc}
\usepackage{textcomp}
\usepackage[utf8]{inputenc}
\usepackage{babel}
\usepackage{CJKutf8}

\usepackage{amsmath, amssymb, amsthm, graphicx}
\usepackage{booktabs, multirow, tabularx}
\usepackage{tikz}
\usetikzlibrary{positioning,arrows.meta,shapes,calc,fit,backgrounds}
\usepackage{pgfplots}
\pgfplotsset{compat=1.16}
\usepackage[most]{tcolorbox}
\tcbuselibrary{skins, breakable}
\usepackage{listings}
\usepackage{xcolor}

\lstset{
  basicstyle=\ttfamily\small,
  keywordstyle=\color{blue!80!black}\bfseries,
  commentstyle=\color{green!50!black}\itshape,
  stringstyle=\color{orange!80!black},
  backgroundcolor=\color{gray!8},
  frame=single,
  rulecolor=\color{gray!40},
  breaklines=true,
  showstringspaces=false,
  numbers=none,
}

\usetheme{metropolis}
\usefonttheme{professionalfonts}
\metroset{sectionpage=none}

\definecolor{LightGray}{RGB}{230,230,230}
\definecolor{RedTitle}{RGB}{0,0,200}
\definecolor{VeryLightGray}{RGB}{245,245,245}
\definecolor{DarkText}{RGB}{0,0,0}
\definecolor{UCRBlue}{RGB}{0,45,106}
\definecolor{UCRGold}{RGB}{241,171,0}

% Stage colours of the course arc (used by the course map and elsewhere)
\colorlet{StageTrunk}{blue!12}
\colorlet{StageBranch}{cyan!14}
\colorlet{StageMachine}{gray!18}
\colorlet{StageWall}{red!14}
\colorlet{StageTools}{orange!20}
\colorlet{StageData}{green!16}
\colorlet{StageAgents}{violet!14}

\hypersetup{bookmarks=false,breaklinks=true,pdfborder={0 0 0},colorlinks=true,
  linkcolor=DarkText,urlcolor=RedTitle}

\setbeamercolor{background canvas}{bg=white}
\setbeamercolor{title}{fg=RedTitle}
\setbeamercolor{frametitle}{bg=VeryLightGray,fg=DarkText}
\setbeamercolor{normal text}{fg=DarkText}
\setbeamercolor{alerted text}{fg=RedTitle}
\setbeamersize{text margin left=5mm, text margin right=5mm}
\addtobeamertemplate{frametitle}{}{\vspace{-0.5em}}

\setbeamertemplate{navigation symbols}{}
\setbeamertemplate{footline}{%
  \hspace{5mm}{\usebeamerfont{page number in head/foot}\color{black!45}%
  ECON 282E \,$\cdot$\, \insertshortdeck}\hfill%
  \usebeamercolor[fg]{page number in head/foot}%
  \usebeamerfont{page number in head/foot}%
  \insertframenumber\,/\,\inserttotalframenumber\kern1em\vskip2pt%
}

% ---------------------------------------------------------------------------
% Chinese text under pdflatex (book titles etc.)
\newcommand{\zh}[1]{\begin{CJK*}{UTF8}{gbsn}#1\end{CJK*}}

% ---------------------------------------------------------------------------
% Deck title block
%   \decktitle{session}{deck id}{title}{date}
\newcommand{\insertshortdeck}{}
\newcommand{\decksession}{}
\newcommand{\deckid}{}
\newcommand{\decktitletext}{}
\newcommand{\deckdate}{}
\newcommand{\decktitle}[4]{%
  \renewcommand{\decksession}{#1}%
  \renewcommand{\deckid}{#2}%
  \renewcommand{\decktitletext}{#3}%
  \renewcommand{\deckdate}{#4}%
  \renewcommand{\insertshortdeck}{Session #1 \,$\cdot$\, Deck #1.#2}%
  \title{#3}%
  \author{Zhigang Feng}%
  \date{#4}%
}
\newcommand{\makedecktitle}{%
  \begin{frame}[plain,noframenumbering]
    \begin{tikzpicture}[remember picture,overlay]
      \fill[UCRBlue] (current page.north west) rectangle ([yshift=-0.55cm]current page.north east);
      \fill[UCRGold] ([yshift=-0.55cm]current page.north west) rectangle ([yshift=-0.62cm]current page.north east);
      \node[anchor=west, text=white, font=\small\bfseries] at ([xshift=5mm,yshift=-0.275cm]current page.north west)
        {ECON 282E \,$\cdot$\, Foundations of Macroeconomics};
      \node[anchor=east, text=white, font=\small] at ([xshift=-5mm,yshift=-0.275cm]current page.north east)
        {UC Riverside \,$\cdot$\, Fall 2026};
    \end{tikzpicture}
    \vspace*{1.1cm}
    {\usebeamerfont{title}\color{black!55}\large Session \decksession\ \,$\cdot$\, Deck \decksession.\deckid\par}
    \vspace{0.25cm}
    {\usebeamerfont{title}\color{RedTitle}\LARGE\bfseries \decktitletext\par}
    \vspace{0.7cm}
    {\large Zhigang Feng\par}
    \vspace{0.15cm}
    {\color{black!65}\deckdate\par}
    \vfill
    {\footnotesize\itshape\color{black!55}Build it on paper $\;\to\;$ solve it on the machine $\;\to\;$ push it to the frontier with AI\par}
  \end{frame}}

% ---------------------------------------------------------------------------
% Section divider (plain frame, not counted)
\newcommand{\sectiondivider}[2]{%
\begin{frame}[plain,noframenumbering]
\begin{center}
\vfill
{\Large\textbf{#1}\par}
\vspace{0.35cm}
{\normalsize #2\par}
\vfill
\end{center}
\end{frame}
}

% ---------------------------------------------------------------------------
% Boxes
\newtcolorbox{conceptbox}[1][]{colback=blue!3, colframe=RedTitle!55, boxrule=0.35mm,
  arc=1mm, left=2mm, right=2mm, top=1mm, bottom=1mm, #1}
\newtcolorbox{cautionbox}[1][]{colback=red!3, colframe=red!50!black, boxrule=0.35mm,
  arc=1mm, left=2mm, right=2mm, top=1mm, bottom=1mm, #1}
\newtcolorbox{examplebox}[1][]{colback=green!3, colframe=green!45!black, boxrule=0.35mm,
  arc=1mm, left=2mm, right=2mm, top=1mm, bottom=1mm, #1}
\newtcolorbox{takeawaybox}[1][]{colback=VeryLightGray, colframe=RedTitle!55, boxrule=0.35mm,
  arc=1mm, left=2mm, right=2mm, top=1mm, bottom=1mm, #1}
\newtcolorbox{researchbox}[1][]{enhanced, colback=violet!4, colframe=violet!60!black,
  boxrule=0.35mm, arc=1mm, left=2mm, right=2mm, top=1mm, bottom=1mm,
  fonttitle=\bfseries\small, title={Research in the wild}, #1}

\newcommand{\compacttables}{\renewcommand{\arraystretch}{1.15}}
\newcommand{\src}[1]{{\tiny\color{black!55}#1}}

% ---------------------------------------------------------------------------
%% Course map: the recurring "you are here" slide for ECON 282E.
%% Loaded by course_preamble.tex. Pattern follows AI-ECON-2026 lec01_map.tex, but the map
%% shows the whole ten-session course rather than one lecture.
%%
%%   \CourseMap{s}{label}   s = 1..10 highlights session s and prints "you are here: label"
%%                          (e.g. \CourseMap{1}{1.A2}); earlier sessions get a check mark.
%%   \CourseMap{0}{}        full overview, nothing highlighted.
%%
%% Note: TikZ option lists are comma-split before TeX conditionals run, so every \ifnum
%% below sits outside [...] option lists.

\newcommand{\CourseMap}[2]{%
\begin{center}
\begin{tikzpicture}[
    sess/.style={rectangle, rounded corners=2pt, draw=black!45, minimum width=1.34cm,
        minimum height=1.12cm, text width=1.26cm, align=center, inner sep=1pt},
    stage/.style={font=\tiny\bfseries, text=black!70},
]
  \foreach \i/\d/\t/\c in {%
      1/Sep 24/Intro \\ Growth/StageTrunk,
      2/Oct 1/HA $\cdot$ OLG \\ Policy/StageBranch,
      3/Oct 8/Num.\ DP \\ Python/StageMachine,
      4/Oct 15/Numerics \\ I/StageMachine,
      5/Oct 22/Numerics \\ II/StageMachine,
      6/Oct 29/The wall \\ \textbf{Midterm}/StageWall,
      7/Nov 5/ML \\ Deep nets/StageTools,
      8/Nov 12/RL \\ HA $+$ DL/StageTools,
      9/Nov 19/Text \\ LLMs/StageData,
      10/Dec 3/Agents \\ \textbf{Projects}/StageAgents} {
    \node[sess, fill=\c] (n\i) at ({1.46*(\i-1)}, 0)
      {{\scriptsize\bfseries S\i}\\[-1pt]{\tiny\color{black!60}\d}\\[1pt]{\tiny \t}};
    \ifnum\i<#1
      \node[font=\scriptsize, text=green!45!black] at ([xshift=-2.5mm,yshift=-2.2mm]n\i.north east) {$\checkmark$};
    \fi
    \ifnum\i=#1
      \draw[RedTitle, very thick, rounded corners=2pt]
        ([xshift=-0.6mm,yshift=-0.6mm]n\i.south west) rectangle ([xshift=0.6mm,yshift=0.6mm]n\i.north east);
      % keep the label inside the picture at both ends of the row
      \ifnum\i<3
        \node[font=\scriptsize\bfseries, text=RedTitle, anchor=south west] at ([yshift=1.2mm]n\i.north west)
          {$\blacktriangledown$\,you are here: #2};
      \else\ifnum\i>8
        \node[font=\scriptsize\bfseries, text=RedTitle, anchor=south east] at ([yshift=1.2mm]n\i.north east)
          {you are here: #2\,$\blacktriangledown$};
      \else
        \node[font=\scriptsize\bfseries, text=RedTitle, anchor=south] at ([yshift=1.2mm]n\i.north)
          {you are here: #2\,$\blacktriangledown$};
      \fi\fi
    \fi
  }
  % stage band under the sessions
  \foreach \a/\b/\lab/\c in {1/1/trunk/StageTrunk, 2/2/branches/StageBranch,
      3/5/the machine $\cdot$ classical methods/StageMachine, 6/6/the wall/StageWall,
      7/8/new tools/StageTools, 9/9/new data/StageData, 10/10/colleagues/StageAgents} {
    \fill[\c] ([yshift=-1.5mm]n\a.south west) rectangle ([yshift=-4.6mm]n\b.south east);
    \path ([yshift=-1.5mm]n\a.south west) -- ([yshift=-4.6mm]n\b.south east)
      node[midway, stage] {\lab};
  }
  \node[font=\footnotesize\itshape, text=black!70] at ([yshift=-9.5mm]$(n1.south)!0.5!(n10.south)$)
    {Build it on paper $\;\to\;$ solve it on the machine $\;\to\;$ push it to the frontier with AI};
\end{tikzpicture}
\end{center}
}


\graphicspath{{./figures/}}

\lstset{language=Python, basicstyle=\ttfamily\scriptsize, frame=none,
        backgroundcolor=\color{gray!8}, aboveskip=2pt, belowskip=2pt}

\decktitle{5}{A1}{What Perturbation Computes}{Thursday, October 22, 2026}

\begin{document}

\makedecktitle

%=============================================================================
\begin{frame}{Where We Are}
\CourseMap{5}{5.A1}
\vspace{-1mm}
\begin{center}
\small \textbf{Block A:} \textbf{5.A1} what perturbation computes $\;\cdot\;$ \textbf{5.A2} the model and the parameter $\;\cdot\;$ \textbf{5.A3} first order $\;\cdot\;$ \textbf{5.A4} linear RE systems $\;\cdot\;$ \textbf{5.A5} higher order, risk, pruning $\;\cdot\;$ \textbf{5.A6} the worked example.
\end{center}
\end{frame}

%=============================================================================
\begin{frame}{Hook: Four Numbers Against Two Thousand}
\begin{columns}[T]
\begin{column}{0.53\textwidth}
\small
Session 4 finished by solving the quarterly RBC model as well as a grid can: continuous choice, linear interpolation, golden section, Howard. It stores \textbf{2{,}100 numbers} --- one per $(k,z)$ --- and takes about a second.
\medskip

By the end of this block we will solve the same model with \textbf{four numbers}, in well under a millisecond, and near the steady state it will be \textbf{more accurate}.
\medskip

That should be alarming. Understanding exactly why it happens --- and exactly where it stops being true --- is the whole of Block A.
\end{column}
\begin{column}{0.44\textwidth}
\begin{center}\footnotesize
\begin{tabularx}{\linewidth}{@{}>{\raggedright\arraybackslash}p{2.35cm} r r@{}}
\toprule
 & \textbf{grid (4.B5)} & \textbf{1st order} \\
\midrule
numbers stored & $2{,}100$   & $\mathbf{4}$ \\
solve time     & $0.94$\,s   & $<0.001$\,s \\
$\max|\mathcal{E}|$ near $k^{*}$ & $10^{-2.82}$ & $\mathbf{10^{-4.01}}$ \\
$\max|\mathcal{E}|$ wide band    & $10^{-2.51}$ & $10^{-2.56}$ \\
\bottomrule
\end{tabularx}
\end{center}
\medskip
\begin{cautionbox}\footnotesize
Read the last row before celebrating. Away from the steady state the advantage is gone.
\end{cautionbox}
\end{column}
\end{columns}
\vspace{1mm}
\src{Ledger \texttt{three\_methods}. ``Near'' is $k\in[0.95k^{*},1.05k^{*}]$, ``wide'' is $[0.5k^{*},1.5k^{*}]$; $\mathcal{E}$ is the unit-free Euler residual of 3.A1.}
\end{frame}

%=============================================================================
\begin{frame}{What We Are Actually Solving}
\begin{columns}[T]
\begin{column}{0.54\textwidth}
\small
Every model in this course reduces to one line: find an unknown \textbf{function} $d$ --- a decision rule --- with
\[
  H(d,\,x,\,\theta) \;=\; 0 ,
\]
where $x$ is the state and $\theta$ the parameters.

Sessions 3 and 4 stored $d$ at $N$ grid points and iterated. That is one way to represent a function.

Perturbation uses a \textbf{series around a point where the answer is known},
\[
  d^{n}(x,\theta)\;=\;\sum_{i=0}^{n} a_i\,(x-x_0)^{i},
\]
with $a_0,\dots,a_n$ chosen so that $H=0$ holds \emph{to order $n$} at $x_0$.
\end{column}
\begin{column}{0.43\textwidth}
\begin{conceptbox}[title=\textbf{The shift}]\small
The unknown stops being \emph{a table of values} and becomes \emph{a short list of derivatives at one point}.
\end{conceptbox}
\medskip
\begin{examplebox}\footnotesize
The known point is the \textbf{deterministic steady state}: computed since 1.B1 in one line, used so far only as a sanity check. It is about to become the whole solution.
\end{examplebox}
\end{column}
\end{columns}
\vspace{1mm}
\src{Quant\_Macro Lec.~9, ``Solving Functional Equations''. Block A's exposition follows Fern\'andez-Villaverde \& Guerr\'on-Quintana, \emph{Solution Methods} V--VI, rewritten here.}
\end{frame}

%=============================================================================
\begin{frame}{The Toolkit: Five Things, and We Have Four of Them}
\begin{columns}[T]
\begin{column}{0.56\textwidth}
\small
\begin{enumerate}\setlength{\itemsep}{1.5mm}
  \item \textbf{Implicit function theorem} --- licenses writing $d$ as a function of $(x,\theta)$ near a reference point at all.
  \item \textbf{Taylor's theorem} --- expands that function, and turns ``solve for $d$'' into ``solve for its derivatives''.
  \item \textbf{Differentiating an identity} --- if $F(\cdot)\equiv 0$ then \emph{every} derivative of $F$ is zero. This is where the equations come from.
  \item \textbf{An eigenvalue condition} --- Blanchard--Kahn, which decides whether the linear system we end up with has one solution, many, or none. \emph{New in 5.A4.}
  \item \textbf{Complementarity reformulation} --- Fischer--Burmeister, when the model has inequalities. \emph{2.A1 and 4.A2.}
\end{enumerate}
\end{column}
\begin{column}{0.41\textwidth}
\begin{takeawaybox}\footnotesize
Only item 4 is genuinely new. Items 1--3 are calculus; item 5 you have already used twice.
\end{takeawaybox}
\medskip
\begin{cautionbox}\footnotesize
Everything here is \textbf{local}. It describes the economy for small deviations from the expansion point and makes no claim beyond them.
\end{cautionbox}
\end{column}
\end{columns}
\vspace{1mm}
\src{Quant\_Macro Lec.~9, ``Perturbation: Mathematical Toolkit''.}
\end{frame}

%=============================================================================
\begin{frame}{The Roadmap, Once, So the Next Three Decks Have a Shape}
\small
\begin{center}
\begin{tabularx}{0.95\linewidth}{@{}p{0.5cm} >{\raggedright\arraybackslash}X p{1.9cm}@{}}
\toprule
 & \textbf{Step} & \textbf{Where} \\
\midrule
1 & Write the equilibrium conditions as one system $F(\cdot\,;\chi)=0$ & 5.A2 \\
2 & Choose the expansion point: deterministic steady state, $z=0$, $\chi=0$ & 5.A2 \\
3 & Replace the policy functions by Taylor series around it & 5.A2 \\
4 & Differentiate $F=0$ there; each derivative is an equation & 5.A3 \\
5 & Solve the linear systems --- first order, then higher orders & 5.A3, 5.A5 \\
6 & Check existence and uniqueness, and ask how far it is valid & 5.A4, 5.A6 \\
\bottomrule
\end{tabularx}
\end{center}
\vspace{1mm}
\begin{center}
\begin{minipage}{0.86\linewidth}
\begin{conceptbox}[title=\textbf{Notice what is missing}]\footnotesize
No grid, no fixed point, no maximization. Steps 4 and 5 are differentiation and linear algebra --- which is the entire reason it is fast.
\end{conceptbox}
\end{minipage}
\end{center}
\vspace{1mm}
\src{Quant\_Macro Lec.~9, ``Roadmap: Solving a DSGE Model by Perturbation''. $\chi$ is the shock-scale parameter of 5.A2.}
\end{frame}

%=============================================================================
\begin{frame}{The Implicit Function Theorem: Permission to Start}
\begin{columns}[T]
\begin{column}{0.55\textwidth}
\small
Before expanding $d$ in a series we need to know that a function $d$ exists at all.
\medskip

\textbf{IFT.} If $H(d_0,x_0,\theta_0)=0$ and
\[
  \frac{\partial H}{\partial d}(d_0,x_0,\theta_0)\;\neq\;0 ,
\]
then on a neighbourhood of $(x_0,\theta_0)$ there is a unique
\[
  d \;=\; f(x,\theta)
\]
solving $H=0$, and it is as smooth as $H$ is.
\medskip

So the non-singularity of one Jacobian buys three things at once: \textbf{existence}, \textbf{local uniqueness}, and \textbf{differentiability} --- and the third is what we are about to spend.
\end{column}
\begin{column}{0.42\textwidth}
\begin{conceptbox}[title=\textbf{Local, and only local}]\small
The theorem says nothing outside the neighbourhood, and never says how large it is. Every limitation of perturbation is already visible in this sentence.
\end{conceptbox}
\medskip
\begin{examplebox}\footnotesize
Compare 3.A1: there, ``solve'' meant a contraction mapping with a global fixed point. Here it means a local inverse. Two different senses of the same word.
\end{examplebox}
\end{column}
\end{columns}
\vspace{1mm}
\src{Quant\_Macro Lec.~9, ``Role of the Implicit Function Theorem''.}
\end{frame}

%=============================================================================
\begin{frame}{Two Derivatives That Look Like a Contradiction}
\begin{columns}[T]
\begin{column}{0.55\textwidth}
\small
Students get stuck here, so it is worth being explicit. Two derivatives appear, and they point opposite ways.
\medskip

\textbf{One must be non-zero:}
\[
  \frac{\partial H}{\partial d}\;\neq\;0
  \quad\text{at }(d_0,x_0,\theta_0).
\]
This is solvability \emph{in the $d$ direction} --- the IFT hypothesis.
\medskip

\textbf{The other must be zero.} Once $d=f(x,\theta)$ is the solution, define
\[
  G(x,\theta)\;:=\;H\big(f(x,\theta),x,\theta\big)\;\equiv\;0 ,
\]
identically, on the whole neighbourhood. So
\[
  \frac{\partial G}{\partial x}=\frac{\partial G}{\partial \theta}=0 .
\]
\end{column}
\begin{column}{0.42\textwidth}
\begin{conceptbox}[title=\textbf{No contradiction}]\small
They are derivatives of \emph{different functions}. $H$ is a map with $d$ as a free argument; $G$ is what is left after the solution has been substituted in --- and it is the zero function.
\end{conceptbox}
\medskip
\begin{takeawaybox}\footnotesize
The first derivative is a \emph{hypothesis}. The second is the \emph{source of every equation} we are going to write down.
\end{takeawaybox}
\end{column}
\end{columns}
\vspace{1mm}
\src{Quant\_Macro Lec.~9, ``Why the Derivatives Matter''.}
\end{frame}

%=============================================================================
\begin{frame}{Is the Hypothesis Satisfied? In the Growth Model, Yes}
\begin{columns}[T]
\begin{column}{0.53\textwidth}
\small
Here $H$ is the pair (Euler equation, resource constraint) and $d=(c,k')$. The relevant Jacobian is
\[
\frac{\partial H}{\partial (c,k')}=
\begin{pmatrix}
\partial_c(\text{Euler}) & \partial_{k'}(\text{Euler}) \\[2pt]
\partial_c(\text{budget}) & \partial_{k'}(\text{budget})
\end{pmatrix}.
\]
It has full rank at the steady state for two economic reasons:
\medskip
\begin{itemize}\small
  \item $u''(c)<0$: concavity makes the Euler equation respond to $c$ at all.
  \item The budget constraint is \textbf{linear} in $(c,k')$, coefficient $1$ on each.
\end{itemize}
\medskip
Were it singular, the model would be degenerate --- one more unit of saving would cost nothing.
\end{column}
\begin{column}{0.44\textwidth}
\begin{conceptbox}[title=\textbf{Where it returns}]\footnotesize
Blanchard--Kahn (5.A4) checks the eigenvalues of the \emph{linearized} system. Passing it confirms, after the fact, that this Jacobian was invertible --- the same condition from the other end.
\end{conceptbox}
\medskip
\begin{examplebox}\footnotesize
$u''<0$ has now done four jobs: uniqueness of the maximum (1.B2), unimodality for golden section (4.A2), the constant in Santos's bound (3.A1), and now local solvability.
\end{examplebox}
\end{column}
\end{columns}
\vspace{1mm}
\src{Quant\_Macro Lec.~9, ``How Is $\partial H/\partial d \neq 0$ Satisfied in Growth Models?''.}
\end{frame}

%=============================================================================
\begin{frame}{The Chain Rule Is the Whole Method}
\begin{columns}[T]
\begin{column}{0.58\textwidth}
\small
With $G(x,\theta)=H\big(f(x,\theta),x,\theta\big)\equiv 0$, differentiate:
\begin{align*}
  \frac{\partial G}{\partial x}
  &= \underbrace{\frac{\partial H}{\partial d}}_{\text{known}}\;
     \underbrace{\frac{\partial f}{\partial x}}_{\textbf{unknown}}
   + \underbrace{\frac{\partial H}{\partial x}}_{\text{known}} \;=\; 0, \\[2mm]
  \frac{\partial G}{\partial \theta}
  &= \frac{\partial H}{\partial d}\;\frac{\partial f}{\partial \theta}
   + \frac{\partial H}{\partial \theta} \;=\; 0 .
\end{align*}
Everything except $\partial f/\partial x$ can be evaluated at the steady state, because we know the steady state. So this is a \textbf{linear equation in the unknown derivative}, and the IFT hypothesis is exactly what makes it solvable.
\medskip

Repeat at second order and you get a linear system for the second derivatives, in which the first derivatives now appear as \emph{known} coefficients. The procedure is recursive.
\end{column}
\begin{column}{0.39\textwidth}
\begin{takeawaybox}\footnotesize
\textbf{Match coefficients, order by order.} Each order is linear given the ones below it --- which is why fifth-order perturbation is not appreciably harder than second, only messier.
\end{takeawaybox}
\medskip
\begin{cautionbox}\footnotesize
The first order is the only one that requires a decision (5.A4). Everything above it is bookkeeping.
\end{cautionbox}
\end{column}
\end{columns}
\vspace{1mm}
\src{Quant\_Macro Lec.~9, ``Chain Rule and Matching Conditions''.}
\end{frame}

%=============================================================================
\begin{frame}{What If the Model Has Inequalities?}
\begin{columns}[T]
\begin{column}{0.54\textwidth}
\small
The method above needs a system of \textbf{equalities}, $F(\cdot)=0$. Models with a borrowing limit or irreversible investment give Kuhn--Tucker conditions instead:
\[
  a\ge 0,\qquad h\ge 0,\qquad a\,h=0 .
\]
The \textbf{Fischer--Burmeister} function turns the triple into one equation:
\[
  \Psi^{FB}(a,h)\;=\;a+h-\sqrt{a^{2}+h^{2}}\;=\;0 .
\]
Check the three cases: $(a>0,h=0)$ gives $a-a=0$; $(a=0,h>0)$ gives $h-h=0$; $(0,0)$ gives $0$. And if both are strictly positive, $\Psi^{FB}<0$ by Cauchy--Schwarz --- so it is zero exactly on the complementarity set.
\end{column}
\begin{column}{0.43\textwidth}
\begin{cautionbox}[title=\textbf{But be honest about it}]\small
$\Psi^{FB}$ is \textbf{not differentiable at the origin} --- precisely where the constraint starts to bind. Perturbation needs derivatives. Reformulating the kink does not remove it.
\end{cautionbox}
\medskip
\begin{examplebox}\footnotesize
This is why Block B exists, and why the Aiyagari economy of 2.A cannot be perturbed at the constraint. We come back to it in 5.A6.
\end{examplebox}
\end{column}
\end{columns}
\vspace{1mm}
\src{Quant\_Macro Lec.~9, ``Perturbation Methods'' (the complementarity frame); same notation as 2.A1 and 4.A2.}
\end{frame}

%=============================================================================
\begin{frame}{Takeaways}
\begin{takeawaybox}
\begin{itemize}\small
  \item Perturbation changes the \textbf{representation} of the unknown: from $N$ values on a grid to a handful of \textbf{derivatives at one point}.
  \item The implicit function theorem is what licenses it, and everything local about the method is already contained in that theorem's fine print.
  \item Two derivatives, opposite requirements: $\partial H/\partial d \neq 0$ is the \emph{hypothesis}; $\partial G/\partial(x,\theta)=0$ is the \emph{source of the equations}.
  \item ``Differentiate an identity, match coefficients, solve a linear system, repeat'' --- and each order is linear given the ones below it.
  \item Inequality constraints can be written as equalities, but the \textbf{kink survives the rewriting}. That limitation is structural, not technical.
  \item Measured, on Session 4's own model: \textbf{four numbers} against $2{,}100$, and better accuracy near the steady state --- but no better far from it.
\end{itemize}
\end{takeawaybox}
\end{frame}

%=============================================================================
\begin{frame}{Bridge: A Model, and One Extra Parameter}
\begin{columns}[T]
\begin{column}{0.53\textwidth}
\small
So far this has been calculus with no economics in it. The next deck puts a model underneath: the stochastic growth model with log utility and full depreciation --- the Brock--Mirman benchmark this course has used since 1.B1, and the one model whose policy we know exactly.
\medskip

It also introduces the trick that makes uncertainty tractable: an artificial parameter that \textbf{scales the shock}, so that ``no risk'' is a point we can expand around.
\end{column}
\begin{column}{0.44\textwidth}
\begin{conceptbox}[title=\textbf{Next (5.A2)}]\small
The model, its closed-form solution, the steady state, and the perturbation parameter $\chi$: why we expand around $\chi=0$ rather than around the model we actually want.
\end{conceptbox}
\medskip
\begin{examplebox}\footnotesize
\textbf{Homework 1} is due today. The practice midterm goes out this afternoon; the midterm itself is next week, Block B, and covers Sessions 1--5.
\end{examplebox}
\end{column}
\end{columns}
\end{frame}

\end{document}
